\clearpage\fancyhead[L]{\small EXCERPT FROM A LETTER.}\hypertarget{native-323-lower}{}\label{native-323-lower}

\begin{center}\textbf{EXCERPT FROM A LETTER BY PROFESSOR GAUSS.}\par\vspace{4pt}\hrule\vspace{4pt}\fontsize{10}{12}\selectfont Monatliche Correspondenz. Vol. XVIII. pp. 289. 290. September 1810.\par\vspace{4pt}\hrule\vspace{4pt}\end{center}
\begin{flushright}Göttingen, 17 August 1810.\end{flushright}

\par\indent ---\enspace---\enspace--- You will already know that our king has now ordered work to continue on our new observatory and with truly royal generosity has granted 200000 francs for it
\SourceLeaf{324}I hope that if you come to see us this autumn, you will find that work has already resumed.

On 8 August I was in Münden, partly to find out how travel affected the rate of our chronometer. The weather, however, was very unfavorable, and between 10 and 11 in the morning I could obtain only a few poor measurements of the Sun's altitude. At noon the Sun was completely hidden, and only after 1 o'clock did I obtain some reliable altitudes. From these, together with some equally good ones taken around 6 o'clock, I found the altitude of the pole to be $51^\circ\ 25'\ 22''$ and the difference in longitude from \emph{Göttingen} to be $1^{\text{m}}9^{\text{s}}3$ west, measured in time. The chronometer performed excellently, and its daily rate from 7\,..\,9 August was $0^{\text{s}}1$, the same as before. Since my solar altitudes were taken so far from noon, I will reserve judgment on the accuracy of my determination until I have a chance to repeat it. I must point out, however, that my nine solar altitudes agree well with one another. My observing site was the so-called \emph{Freitagswerder}, at the confluence of the \emph{Fulda} and the \emph{Werra}, just north of the town. In 1794 Professor \textsc{Seyffer} found the altitude of the pole to be $51^\circ\ 26'\ 52''$ at a spot probably a little farther \emph{south}, though he used only a single altitude.
